\documentclass[11pt,a4paper]{article}

\usepackage[margin=2.4cm,top=2.6cm,bottom=2.6cm]{geometry}
\usepackage{amsmath,amssymb}
\usepackage{booktabs}
\usepackage[table]{xcolor}
\usepackage{listings}
\usepackage{fancyhdr}
\usepackage{titlesec}
\usepackage{enumitem}
\usepackage{parskip}
\usepackage[T1]{fontenc}
\usepackage[expansion=false]{microtype}
\usepackage[hidelinks,colorlinks=true,linkcolor=accent,urlcolor=accent]{hyperref}

\definecolor{accent}{HTML}{1F4E79}
\definecolor{rule}{HTML}{B8C4CE}
\definecolor{codebg}{HTML}{F6F7F9}
\definecolor{warn}{HTML}{9A4B00}

\titleformat{\section}{\Large\bfseries\color{accent}}{\thesection}{0.7em}{}
\titleformat{\subsection}{\large\bfseries}{\thesubsection}{0.6em}{}
\titlespacing*{\section}{0pt}{1.6em}{0.7em}

\pagestyle{fancy}\fancyhf{}
\renewcommand{\headrulewidth}{0.4pt}
\fancyhead[L]{\small\textcolor{accent}{creditlib \textbar\ Endogenous Recovery}}
\fancyhead[R]{\small\textcolor{accent}{Design \& Process}}
\fancyfoot[C]{\small\thepage}

\lstset{basicstyle=\ttfamily\scriptsize,backgroundcolor=\color{codebg},
  frame=single,rulecolor=\color{rule},breaklines=true,showstringspaces=false,
  language=Python,xleftmargin=0pt,framexleftmargin=3pt,aboveskip=0.9em,belowskip=0.9em}

\newcommand{\keypoint}[1]{\par\vspace{0.4em}\noindent
  \colorbox{codebg}{\parbox{\dimexpr\linewidth-2\fboxsep}{\small #1}}\par\vspace{0.5em}}

\begin{document}

\begin{titlepage}
\centering
\vspace*{2.2cm}
{\Huge\bfseries\color{accent} Endogenous Recovery\par}
\vspace{0.7em}
{\LARGE Design and Process Specification\par}
\vspace{0.5em}
{\large Extending \texttt{creditlib} beyond the flat 40\% convention\par}
\vspace{2.2cm}
\rule{0.62\textwidth}{0.6pt}\par
\vspace{1.1em}
{\large Shreyas Sundaresan\par}
\vspace{0.4em}
{\normalsize Design specification, methodology, and implementation plan\par}
\vspace{1.1em}
\rule{0.62\textwidth}{0.6pt}\par
\vspace{2.2cm}

\begin{minipage}{0.86\textwidth}\small
\textbf{Purpose.} Replace the market-convention assumption of a uniform 40\%
recovery rate with issuer-specific estimates produced by a fundamentals-based model,
and propagate those estimates through the CDS bootstrap so that implied default
probabilities reflect loss given default rather than assuming it away.

\medskip
\textbf{Central constraint.} A CDS spread identifies expected loss, approximately
$h(1-R)$ --- one equation in two unknowns. Recovery is therefore not identifiable from
spreads alone. This document specifies how that constraint is broken, why the chosen
route avoids circularity, and what changes analytically as a result.

\medskip
\textbf{Status.} Model implemented and integrated. Coefficients are calibrated to
published anchors, \emph{not} estimated from a default database. 104 tests passing.
\end{minipage}

\vfill
{\small Version 1.0 \textbar\ August 2026}
\end{titlepage}

\tableofcontents
\newpage

%% ============================================================= 1
\section{Objective and Scope}

\subsection{What this extends}

The existing \texttt{creditlib} engine bootstraps a piecewise-flat hazard curve from
quoted CDS spreads, prices both legs by exact closed-form integration, and computes
standard credit risk measures. It assumes a single recovery rate --- conventionally
40\% --- applied uniformly to every issuer.

This extension replaces that assumption with a per-issuer estimate derived from
fundamentals, and threads it through the bootstrap.

\subsection{Two requirements, one of which is trivial}

\textbf{Scale to twenty or more issuers.} Already satisfied. The library imposes no
limit on the number of names; a forty-name cross-section bootstraps in under four
seconds. No development required, and none was undertaken.

\textbf{Endogenise recovery.} The substantive work, and the subject of this document.

\subsection{Why the flat convention is not merely imprecise}

A uniform 40\% is defensible as a \emph{quoting convention}: it makes spreads
comparable across issuers by fixing the second unknown. It is not defensible as an
estimate, and it is materially wrong for identifiable categories.

The clearest case is bank holding company debt. Single-name bank CDS references the
\emph{holding company}, and under Dodd-Frank Title II with a single-point-of-entry
resolution strategy, holdco senior unsecured is explicitly TLAC --- it exists to be
written down or converted so that operating subsidiaries continue. Its recovery in
the state that triggers a credit event is designed to be low. Lehman Brothers senior
settled at 8.625 cents in the October 2008 auction, roughly one fifth of the convention.

\keypoint{Applying 40\% to a bank holdco and to a secured real estate credit treats
two structurally different claims as identical. The resulting default probabilities
are not comparable, and the error is systematic rather than random.}

%% ============================================================= 2
\section{The Identification Problem}

\subsection{Statement}

To first order, a CDS par spread satisfies the credit triangle
\[
s \;\approx\; h\,(1-R)
\]
and exactly, in the continuous-premium limit with flat rates, $s = h(1-R)$.

One observable, two unknowns. The pair $(h, R)$ is not identified: for any assumed
$R$ there exists an $h$ reproducing the observed spread exactly.

\subsection{Demonstration}

Bootstrapping an identical five-tenor quote set at $R = 40\%$ and $R = 20\%$ produces:

\begin{center}
\begin{tabular}{lrr}
\toprule
& \textbf{$R = 40\%$} & \textbf{$R = 20\%$} \\
\midrule
Repricing error (relative) & $4.4\times10^{-16}$ & $4.4\times10^{-16}$ \\
Hazard ratio to the 40\% case & 1.000 & 0.750 \\
Implied 5Y default probability & 8.37\% & 6.34\% \\
Implied ranking across names & unchanged & unchanged \\
\bottomrule
\end{tabular}
\end{center}

Both fit perfectly. The hazards differ by exactly $(1-0.40)/(1-0.20) = 0.75$. This is
not a numerical artefact --- it is the identification failure made visible, and it is
covered by a regression test.

\subsection{The one thing the degeneracy protects}

Under a \emph{uniform} recovery, $h = s/(1-R)$ is a strictly monotone transform of the
spread. The implied default probability ranking is therefore invariant to the
assumption. Levels move; order does not.

This is why the existing engine's conclusions survive the weak assumption. It is also
precisely what the extension gives up.

%% ============================================================= 3
\section{Routes to Identification}

Three ways to obtain a second equation. All were considered.

\subsection{Route A --- Fundamentals model (selected)}

Estimate $R$ from issuer characteristics outside the CDS market entirely: seniority,
sector asset tangibility, debt cushion, leverage, macroeconomic state. Feed the result
into the bootstrap as a known input.

\textbf{Advantages.} No additional market data required. No circularity --- nothing in
the recovery model reads a spread. Applies to every issuer, including illiquid names
with a single quote. Directly interpretable, and the drivers are the same ones a
fundamental credit analyst already examines.

\textbf{Disadvantages.} The estimate is only as good as the model. Requires a default
database for genuine statistical estimation; without one it is an expert scorecard.
Introduces model risk where previously there was only an acknowledged convention.

\subsection{Route B --- Market-implied via a second instrument}

Recovery becomes identifiable given a second quote on the same issuer with different
loss given default. Practical instruments:

\begin{itemize}[leftmargin=1.4em,itemsep=0.15em]
\item \textbf{Senior versus subordinated CDS.} Same hazard, different recovery. Two
equations, two unknowns.
\item \textbf{Recovery swaps or fixed-recovery CDS.} Directly quoted recovery.
\item \textbf{CDS--bond basis.} The cash bond price contains recovery information.
\end{itemize}

\textbf{Advantages.} Market-implied rather than modelled; reflects the market's own view.

\textbf{Disadvantages.} Requires two liquid quotes per issuer. Subordinated CDS is thin
for most names and effectively absent for many. Recovery swaps are rare. Restricts the
universe to the most liquid handful of issuers, which defeats the purpose of a
twenty-name cross-section.

\subsection{Route C --- Bayesian combination}

Fundamentals model as prior; market data as likelihood where a second instrument exists.

\textbf{Assessment.} The correct long-run architecture, and the natural version~2. Not
attempted here: without a fitted prior it would layer complexity on an unestimated model.

\subsection{Decision}

\keypoint{\textbf{Route A is implemented.} Route B is specified as the version~2
extension for the subset of issuers with a senior/subordinated pair, at which point
Route C becomes available for free. The architecture accepts a per-name recovery
mapping regardless of its provenance, so switching or blending routes requires no
change to the pricing path.}

%% ============================================================= 4
\section{Model Specification}

\subsection{Functional form}

A logit-linear scorecard:
\[
z \;=\; \beta_0 \;+\; \beta_{\text{sen}}[\text{seniority}]
\;+\; \beta_{\text{tang}}(\tau - 0.5)
\;+\; \beta_{\text{cush}}\,c
\;-\; \beta_{\text{lev}}\frac{\ell - 3}{10}
\;+\; \beta_{\text{macro}}\,m
\]
\[
R \;=\; \frac{1}{1 + e^{-z}}
\]

The logistic link is not cosmetic. The bootstrap divides by $(1-R)$, so a prediction of
exactly $1.0$ would be a division by zero and a prediction below $0$ would be
meaningless. The link guarantees $R \in (0,1)$ for \emph{any} input combination, which
is enforced by a test sweeping extreme values of every covariate.

\subsection{Covariates and rationale}

\begin{center}
\begin{tabular}{p{2.6cm}p{3.0cm}p{8.4cm}}
\toprule
\textbf{Covariate} & \textbf{Measure} & \textbf{Rationale} \\
\midrule
Seniority & Categorical & Position in the capital structure is the dominant driver.
Secured claims recover most; junior subordinated least. \\[0.3em]
Tangibility $\tau$ & Sector proxy, $[0,1]$ & Hard assets survive bankruptcy and can be
sold. Utilities and real estate recover far more than software. \\[0.3em]
Debt cushion $c$ & Fraction of debt ranking junior & Junior debt absorbs losses first,
lifting recovery on the senior claim. \\[0.3em]
Leverage $\ell$ & Net debt / EBITDA & Less asset value per unit of claim. Neutral at 3.0. \\[0.3em]
Macro state $m$ & $[-1,1]$ & Recoveries fall in downturns. Recovery and default are
positively correlated, which is exactly when it hurts. \\
\bottomrule
\end{tabular}
\end{center}

\subsection{Calibration anchors}

Coefficients are set so that neutral covariates reproduce the published ordering:

\begin{center}
\begin{tabular}{lrr}
\toprule
\textbf{Seniority} & \textbf{Target} & \textbf{Model} \\
\midrule
Secured & 65\% & 65.6\% \\
Senior unsecured & 40\% & 40.0\% \\
Bank holdco senior (TLAC) & 24\% & 24.0\% \\
Subordinated & 20\% & 18.9\% \\
Junior subordinated & 10\% & 9.5\% \\
\bottomrule
\end{tabular}
\end{center}

\subsection{Status: calibrated, not estimated}

\keypoint{\textbf{This is an expert scorecard with a defensible functional form. It is
not a fitted statistical model.} Coefficients are calibrated to published anchors
rather than estimated from a recovery database. Any write-up, presentation or interview
discussion must describe it as calibrated. Overstating it is the fastest available way
to lose credibility on this work, and the distinction is one an experienced credit
analyst will probe within two questions.}

The estimation path exists. \texttt{RecoveryModel.fit()} performs ridge-regularised
least squares in logit space, requires a minimum of twelve observations, and sets the
uncertainty band from residual dispersion. It refuses to run on smaller samples on the
grounds that the calibrated defaults are more defensible than an overfitted estimate.
The \texttt{fitted} flag stays \texttt{False} until real data is supplied.

\subsection{Uncertainty band}

The model reports an interval, not a point. The default half-width of 0.55 logit units
gives roughly $\pm 12$ points around a 40\% central estimate. This is deliberately wide:
empirical recovery distributions are broad and materially bimodal, and a narrow band
would be false precision. After \texttt{fit()}, the band is replaced by the residual
standard deviation.

%% ============================================================= 5
\section{Architecture and Integration}

\subsection{Separation of concerns}

\begin{lstlisting}[language={}]
creditlib/
  recovery.py       IssuerProfile -> RecoveryModel -> RecoveryEstimate
  analyze.py        compare(names, recovery=float | {name: float})
  bootstrap.py      unchanged -- already accepted a recovery argument
  cds.py            unchanged
\end{lstlisting}

\keypoint{\textbf{The pricing path is untouched.} \texttt{bootstrap\_survival\_curve}
already took recovery as an argument; the extension changes only what is passed to it.
That is the test of a correct architecture --- adding a recovery model required no
change to any pricing code, and all 86 pre-existing tests continued to pass unmodified.}

\subsection{Interface}

\begin{lstlisting}
from creditlib import RecoveryModel, IssuerProfile, recovery_map, compare

profiles = [
    IssuerProfile("JPMorgan",   sector="banks",     is_bank=True, debt_cushion=0.22),
    IssuerProfile("Duke Energy", sector="utilities", leverage=4.8, debt_cushion=0.05),
    IssuerProfile("Prologis",   sector="real_estate", seniority="secured",
                  leverage=5.2, debt_cushion=0.15),
]

model = RecoveryModel()
R = recovery_map(model, profiles)          # {name: recovery}

compare({"JPMorgan": 58, "Duke Energy": 62, "Prologis": 57}, recovery=R)
\end{lstlisting}

Recovery may be a scalar (uniform, as before) or a mapping (per-name). Supplying a
mapping with a missing name raises rather than silently defaulting --- a silent fallback
to 40\% would reintroduce the exact error the module exists to remove.

\subsection{Automatic bank reclassification}

Constructing an \texttt{IssuerProfile} with \texttt{is\_bank=True} and seniority
\texttt{senior\_unsecured} silently promotes it to \texttt{holdco\_senior} and records
the reason in the profile's notes. Single-name bank CDS references the holdco, and
treating that claim as ordinary corporate senior is the largest single error the module
is designed to prevent. Making it an automatic correction rather than user
responsibility is a deliberate choice.

%% ============================================================= 6
\section{Empirical Results}

\subsection{Test universe}

Twenty issuers across ten sectors and two seniority classes: six financials, two
utilities, two energy, one transport, one autos, one consumer staples, one retail, one
healthcare, one pharma, one software, one media, one real estate (secured), one telecom.
Spreads span 41bp to 395bp.

\subsection{Predicted recovery dispersion}

\begin{center}
\begin{tabular}{lrl}
\toprule
\textbf{Issuer} & \textbf{Predicted $R$} & \textbf{Principal driver} \\
\midrule
Prologis (secured, real estate) & 74.6\% & Secured claim, high tangibility \\
NextEra (utilities) & 51.8\% & Hard assets \\
Exxon Mobil (energy) & 51.2\% & Hard assets, low leverage \\
Kroger (consumer staples) & 42.7\% & Near the convention \\
Oracle (software) & 31.6\% & Low tangibility \\
Citigroup (bank holdco) & 24.2\% & TLAC layer \\
U.S. Bancorp (bank holdco) & 21.7\% & TLAC, thin debt cushion \\
\bottomrule
\end{tabular}
\end{center}

Range: 21.7\% to 74.6\%. The flat convention assigns 40\% to all twenty.

\subsection{Effect on the ranking}

Recomputing implied default probabilities with model recovery in place of the flat 40\%:

\begin{center}
\begin{tabular}{lrrrr}
\toprule
\textbf{Issuer} & \textbf{Spread} & \textbf{PD (flat 40\%)} & \textbf{PD (model)} & \textbf{Rank move} \\
\midrule
JPMorgan & 58 bp & 4.76\% & 3.76\% & $+4$ \\
Bank of America & 63 bp & 5.16\% & 4.06\% & $+5$ \\
Wells Fargo & 72 bp & 5.87\% & 4.60\% & $+5$ \\
Citigroup & 85 bp & 6.90\% & 5.50\% & $+5$ \\
Exxon Mobil & 41 bp & 3.39\% & 4.15\% & $-3$ \\
NextEra & 55 bp & 4.52\% & 5.59\% & $-5$ \\
Prologis & 57 bp & 4.68\% & 10.72\% & $-10$ \\
\bottomrule
\end{tabular}
\end{center}

\keypoint{\textbf{Thirteen of twenty issuers change rank. The largest move is ten
places.} Prologis --- secured, 74.6\% recovery --- moves from fifth-safest to
fifteenth: because each default costs little, the same 57bp spread requires far more
defaults to justify. Every bank moves the other way, because a low-recovery claim needs
fewer defaults to justify the same premium.}

\subsection{The finding that matters}

The flat convention \textbf{systematically overstates implied default probabilities for
bank holding companies relative to asset-heavy credits}, because it assigns them a
recovery rate roughly sixteen points above what their position in the resolution
hierarchy implies.

This is directly relevant to comparing bank credit against non-financial credit on a
CDS-implied basis. Any cross-sector ranking computed at a flat 40\% carries this bias,
and the bias is one-directional rather than noise.

\subsection{What is now given up}

Ranking invariance. Under a uniform recovery the ordering was independent of the
assumption; it now depends on the recovery model. That is the intended trade: the
ranking carries an analytical view rather than being a monotone restatement of the
spreads. It must be presented as such, with the model's limitations attached.

%% ============================================================= 7
\section{Validation}

\subsection{Implemented (18 new tests, 104 total)}

\begin{itemize}[leftmargin=1.4em,itemsep=0.15em]
\item \textbf{Calibration anchors} reproduced within 2 percentage points.
\item \textbf{Strict seniority ordering} across all five classes.
\item \textbf{Boundedness} --- $R \in (0,1)$ under a sweep of extreme leverage, cycle,
cushion and seniority combinations. Guards the division by $(1-R)$.
\item \textbf{Driver directions} --- cushion up, leverage down, recession down,
tangibility up. Each asserted separately.
\item \textbf{Bank auto-reclassification} to the TLAC layer, with the note recorded.
\item \textbf{Identification property} --- uniform recovery leaves ranking invariant;
name-specific recovery breaks it. Both directions asserted.
\item \textbf{Exact repricing preserved} under name-specific recovery
($< 10^{-13}$ relative).
\item \textbf{Missing recovery rejected}, never silently defaulted.
\item \textbf{Estimator correctness} on a planted synthetic signal, and ridge shrinkage
pinned directionally so it cannot be lost silently.
\end{itemize}

\subsection{Not yet possible}

\keypoint{\textbf{There is no out-of-sample validation, because there is no sample.}
The model is calibrated, not estimated, so accuracy against realised recoveries is
untested. This is the single largest open item and no claim about predictive
performance should be made until it is closed.}

\subsection{Validation plan once data is available}

\begin{enumerate}[leftmargin=1.4em,itemsep=0.2em]
\item \textbf{Assemble} realised recoveries --- Moody's Ultimate Recovery Database, S\&P
LossStats, or credit-event auction settlement prices, which are public and a viable
starting point.
\item \textbf{Fit} via \texttt{RecoveryModel.fit()} with ridge regularisation.
\item \textbf{Cross-validate} by time period, not randomly. Recovery is strongly
cyclical, so random folds leak macro state across the split and flatter the model.
\item \textbf{Benchmark} against two naive baselines: the flat 40\% convention, and a
seniority-only lookup. If the full model does not beat seniority-only, the additional
covariates are not earning their place.
\item \textbf{Report} mean absolute error in recovery points, and separately the effect
on implied default probability, which is what the analysis actually consumes.
\item \textbf{Stability} --- refit annually, monitor coefficient drift.
\end{enumerate}

%% ============================================================= 8
\section{Limitations}

\begin{enumerate}[leftmargin=1.4em,itemsep=0.3em]
\item \textbf{Calibrated, not estimated.} Stated repeatedly because it is the point
most likely to be overstated.

\item \textbf{Deterministic recovery.} The model returns a point estimate with a band,
but the bootstrap consumes only the point. Recovery is genuinely stochastic and
correlated with default. A full treatment requires a joint distribution.

\item \textbf{Sector tangibility is a coarse proxy.} A single number per sector ignores
within-sector dispersion, which is wide.

\item \textbf{No jurisdiction effect.} Bankruptcy regimes differ materially --- US
Chapter 11 versus UK administration versus European insolvency codes. Not modelled.

\item \textbf{Static macro state.} A single scalar, user-supplied. A term structure of
recovery --- lower in the near term during a downturn --- is not represented.

\item \textbf{Bank treatment is structural, not quantitative.} The TLAC reclassification
correctly identifies that holdco senior sits in the bail-in layer, but the resulting
24\% anchor is judgement informed by one observation (Lehman) and the design intent of
the resolution regime. It is not an estimate.
\end{enumerate}

%% ============================================================= 9
\section{Roadmap}

\begin{center}
\begin{tabular}{p{1.1cm}p{5.4cm}p{8.0cm}}
\toprule
\textbf{Ver.} & \textbf{Scope} & \textbf{Unlocks} \\
\midrule
1.0 & Fundamentals scorecard, per-name recovery through the bootstrap & Cross-sector
comparison free of the flat-40\% bias \\[0.3em]
1.1 & Fit on auction settlement data & Moves from calibrated to estimated; enables
out-of-sample error reporting \\[0.3em]
2.0 & Senior/subordinated basis for issuers with both quoted & Market-implied recovery;
identification without a model \\[0.3em]
2.1 & Bayesian combination of model prior and market likelihood & Full coverage with
market updating where available \\[0.3em]
3.0 & Stochastic recovery with default correlation & Removes the deterministic-recovery
limitation; required for portfolio work \\
\bottomrule
\end{tabular}
\end{center}

\subsection{Immediate next step}

Version 1.1. Credit-event auction settlement prices are public, span roughly two decades,
and cover several hundred defaults. That is sufficient to move the headline claim from
``calibrated to published anchors'' to ``estimated from realised recoveries, with
out-of-sample error of $X$ points'' --- a materially stronger statement, and the one
worth making.

\vspace{2em}
\hrule
\vspace{0.8em}
{\small\itshape
Spreads used are indicative and supplied as inputs; the discount curve is illustrative.
Recovery estimates are model output from calibrated coefficients and have not been
validated against realised recoveries. Nothing in this document is investment advice.}

\end{document}
